\subsection{Calibration under standard EPA and EVA profiles}
\label{sec:standard_results}
The standard power-delay profiles reproduce the high-SNR calibration gap
without changing the receiver or its search settings
(Tables~\ref{tab:standard_nmse} and~\ref{tab:standard_gain}). At 20~dB,
the mean residual noise estimates are 14.366 and 11.234 times the thermal
variance for EPA and EVA, respectively. The projected-covariance predictions
are 14.316 and 11.288, whereas the CP estimates have mean ratios 1.003 and
1.001. Thus the externally specified profiles exhibit the same distinction
between averaging a leaked residual and obtaining a signal-cancelling
noise measurement.

For the continuous-delay receiver, CP calibration improves 20~dB NMSE by
5.75~dB on EPA and 6.12~dB on EVA, with paired 95\% intervals
[5.40, 6.07] and [5.91, 6.31]~dB. The corresponding NMSE values move from
$-18.04$ to $-23.79$~dB and from $-13.22$ to $-19.34$~dB; the known-noise
references are $-23.87$ and $-19.36$~dB. Uncoded BER falls from 0.0213 to
0.0086 on EPA and from 0.0288 to 0.0106 on EVA. All three seed batches
give positive high-SNR NMSE gains: 5.46--5.90~dB for EPA and 5.94--6.30~dB
for EVA. Integer EB gains 0.33 and 0.60~dB at the same SNR, retaining its
integer-subspace representation error despite improved calibration.

The SNR sweep locates the useful operating regime. At 10~dB, the
continuous-delay NMSE gains are 0.32~dB for EPA and 0.94~dB for EVA. At
0~dB, they are $-0.17$ and $-0.07$~dB, while BER decreases in both cases.
These metric-dependent low-SNR outcomes accompany a residual inflation of
approximately 10--13\%, compared with more than an order of magnitude at
20~dB. The benefit of changing the noise source therefore follows the
calibration regime rather than constituting a uniform receiver improvement.

\begin{table}[t]
\centering
\caption{Standard EPA/EVA power-delay profiles: noise calibration and continuous-delay receiver NMSE (dB), with 900 paired groups per condition. R and CP use 16-block residual and CP calibration; K uses known noise. Noise ratios are empirical means divided by the thermal variance. All receiver settings are fixed.}
\label{tab:standard_nmse}
\small
\setlength{\tabcolsep}{5pt}
\begin{tabular}{llrrrrr}
\toprule
PDP & SNR & \multicolumn{2}{c}{Noise ratio} & \multicolumn{3}{c}{NMSE (dB)} \\
\cmidrule(lr){3-4}\cmidrule(lr){5-7}
 & (dB) & R & CP & R & CP & K \\
\midrule
EPA & 0 & 1.126 & 0.998 & -6.91 & -6.73 & -6.78 \\
EPA & 10 & 2.346 & 1.004 & -14.61 & -14.92 & -15.03 \\
EPA & 20 & 14.366 & 1.003 & -18.04 & -23.79 & -23.87 \\
\midrule
EVA & 0 & 1.104 & 1.002 & -3.60 & -3.53 & -3.57 \\
EVA & 10 & 2.031 & 1.001 & -10.24 & -11.18 & -11.18 \\
EVA & 20 & 11.234 & 1.001 & -13.22 & -19.34 & -19.36 \\
\bottomrule
\end{tabular}
\end{table}

\begin{table}[t]
\centering
\caption{Paired NMSE gains from CP calibration and uncoded BER for the same standard-profile experiment. Positive gain favors CP; brackets give pointwise 95\% paired bootstrap intervals, stratified by seed. BER columns use the continuous-delay receiver with residual (R), CP, and known-noise (K) calibration.}
\label{tab:standard_gain}
\small
\setlength{\tabcolsep}{5pt}
\begin{tabular}{llrrrr}
\toprule
PDP & SNR & NMSE gain (dB) & \multicolumn{3}{c}{BER} \\
\cmidrule(lr){4-6}
 & (dB) & CP over R [95\% CI] & R & CP & K \\
\midrule
EPA & 0 & -0.17 [-0.33, -0.02] & 0.2686 & 0.2627 & 0.2632 \\
EPA & 10 & 0.32 [0.08, 0.57] & 0.0634 & 0.0566 & 0.0570 \\
EPA & 20 & 5.75 [5.40, 6.07] & 0.0213 & 0.0086 & 0.0086 \\
\midrule
EVA & 0 & -0.07 [-0.15, 0.02] & 0.2972 & 0.2937 & 0.2930 \\
EVA & 10 & 0.94 [0.82, 1.06] & 0.0800 & 0.0716 & 0.0706 \\
EVA & 20 & 6.12 [5.91, 6.31] & 0.0288 & 0.0106 & 0.0106 \\
\bottomrule
\end{tabular}
\end{table}


The standard-profile experiment retains all 16\,200 continuous-delay fits,
including 1329 with a nonconverged candidate and 493 with a nonconverged
selected candidate at the 60-evaluation limit. Replaying every flagged fit
with a 300-evaluation limit leaves 92 nonconverged selected candidates.
The six aggregate CP-over-residual NMSE gains change by at most 0.0015~dB;
the two 20~dB gains change by less than 0.001~dB. The primary tables retain
the fixed 60-evaluation budget. This audit establishes stability of the
reported calibration contrast to this budget increase, not global convergence.
